\subsection{The ordinary integral tangent and its isogeny orbit}

\begin{result}{A.4.37}{Integral tangent, isogeny law, and branch denominators}{Proved}
At a place of good ordinary reduction of residue characteristic
$p>3$, the inverse-Verschiebung construction gives an integral line
in $\operatorname{Lie}E^\dagger$ projecting isomorphically to
$\operatorname{Lie}E$. Write its slope as $\upsilon_p(E)$ in the
algebraic coframe $(du,ds)$ of A.4.14 and a good short N\'eron frame.
It is functorial under actual isogenies.

For a finite subgroup $C$ of order $N$, put
$S_C=\sum_{P\in C\setminus\{0\}}x(P)$, and let $f:E\to E'$ be its
quotient, with $f^*\omega'=\lambda\omega$. In short frames,
\[
 \lambda^2h(E')=Nh(E)+S_C,\qquad
 \lambda^2\upsilon_p(E')=N\upsilon_p(E)+S_C.
\]
For $h\in\Qb$, fix an embedding of a number field containing
$h$ and the isogeny data into $\mathbf C_p$. In that field,
\begin{equation}\label{ell:ordinary-isogeny-difference}
 \lambda^2\bigl(h(E')-\upsilon_p(E')\bigr)
       =N\bigl(h(E)-\upsilon_p(E)\bigr).
\end{equation}
For every number field $K$ containing $g_2,g_3,h$, the actual Hodge
coefficients in differential-preserving quotient frames satisfy
$\sigma(h(E_C))=h(E_{\sigma C})$ for $\sigma\in\operatorname{Gal}(\Qb/K)$.

Normalize $v_p(p)=1$. Suppose $h$ is integral at the chosen place and
$h\ne\upsilon_p(E)$, and put $b=v_p(h-\upsilon_p(E))\ge0$.
Among cyclic subgroups of order $p^n$, let $m$ be the logarithm to
base $p$ of the order of the intersection with multiplicative torsion.
Then, in good N\'eron frames on their targets,
\[
 v_p(\lambda_C)=m,\qquad
 d_p(C):=\max\{0,-v_p(h(E_C))\}=\max\{0,2m-n-b\}.
\]
There are respectively
\[
 p^n\ (m=0),\qquad
 (p-1)p^{n-m-1}\ (1\le m<n),\qquad 1\ (m=n)
\]
such subgroups; their total is $D_n=(p+1)p^{n-1}$.
The sum of $d_p(C)$ is at most $S_n$, where
\[
 S_{2r}=\frac{2(p^r-1)}{p-1},\qquad
 S_{2r+1}=\frac{(p+1)p^r-2}{p-1}.
\]
For non-CM $E$, every actual $\operatorname{Gal}(\Qb/K)$-orbit
$\mathcal O$ of these subgroups satisfies
\[
 \frac1{|\mathcal O|}\sum_{C\in\mathcal O}d_p(C)
       \le I_p\frac{S_n}{D_n}=O_{E,K,p}(p^{-n/2}),
\]
where $I_p=[\mathrm{GL}_2(\mathbf Z_p):\rho_p(\operatorname{Gal}(\Qb/K))]$.
\end{result}

\begin{proof}
Work first on the formally \`etale ordinary moduli neighborhood in
\cite[\S2.2, pp.16--19; Theorem 2.2]{EllMochizuki2000Connections}.
Its maps $V^j:E^{F^j}\to E$ are finite \`etale. An identity
first-order point of $E$ therefore has a unique compatible lift
through this tower. The theorem identifies its inverse limit with
the completion of the source's integral universal extension.
Pull the entire tower and this identification along the classifying
map of the given ordinary elliptic scheme; this does not assert that
an individual classifying point is \`etale. On tangent spaces the
projection to $\operatorname{Lie}E$ is an isomorphism, giving the
claimed integral line inside $\operatorname{Lie}E^\dagger$.
The integral short de Rham classes form a basis: the order-two
principal part of $\xi$ is removed by the differential of an
integral simple pole, and its cup pairing with $\omega$ is a unit.
Thus its slope $\upsilon_p(E)$ is integral in the indicated coframe.

The disappearance of a vertical integral tangent can also be read
in the actual binomial algebra of the source's Lemma 2.1, p.16:
\[
 \left.\frac{d}{dT}\binom{T}{r}\right|_{T=0}
       =\frac{(-1)^{r-1}}r.
\]
For a putative vertical derivative $b\partial_T$, degrees $r=p^j$
require $v_p(b)\ge j$ for every $j$, hence $b=0$.

For functoriality, an isogeny $f:E\to E'$ induces actual maps
$f^{F^j}:E^{F^j}\to(E')^{F^j}$ commuting with the two Verschiebung
maps. Modulo $p^j$, compose them with the normalized splittings
$\kappa[p^j]$ in the same source. The two resulting maps to
$(E')^\dagger$ lift the same map to $E'$. Their difference is a map
from a proper elliptic scheme to its vector kernel. Its global
regular functions are those of the base, so this difference is
constant; it is zero at the identity. Passing to inverse limits and
tangents proves the asserted functoriality.

Here is the exact second-kind formula needed for the isogeny law.
For the differential-preserving quotient, its period lattice
$\Lambda_C$ contains $\Lambda$ with index $N$. Comparison of poles
and the constant Laurent term gives
\[
 \wp_{\Lambda_C}(u)=\sum_{v\in C}\wp_\Lambda(u+v)-S_C.
\]
For $v\ne0$, put
$F_v=(y-y(v))/(2(x-x(v)))$; the addition formula gives
$dF_v=(\wp(u)-\wp(u+v))\,du$. Therefore
\[
 f_C^*\xi_C=N\xi-S_C\omega-d\sum_{v\ne0}F_v.
\]
Pullback preserves $H^{0,1}$, so $h(E_C)=Nh(E)+S_C$ in this frame.
Its formulas are algebraic. Thus a $\sigma$ fixing $g_2,g_3,h$
sends this actual equality to the equality for $\sigma C$, proving
the stated $\operatorname{Gal}(\Qb/K)$ compatibility. There is no
claim here for a $\sigma$ not fixing $h$.

For a general short frame with $f^*\omega'=\lambda\omega$, one has
$\omega_C=\omega'/\lambda$, $x_C=\lambda^2x'$ and
$\xi_C=\lambda\xi'$. Hence in algebraic de Rham cohomology,
\[
 f^*\xi'=\frac N\lambda\xi-\frac{S_C}{\lambda}\omega,
 \qquad
 (u,s)\longmapsto
 \left(\lambda u,\frac N\lambda s+\frac{S_C}{\lambda}u\right)
\]
on universal-extension Lie coordinates. Apply this same map to the
Hodge line and to the proved functorial integral line. It gives both
slope identities and \eqref{ell:ordinary-isogeny-difference}.

For the local multiplier, factor $f_C$ through the multiplicative
intersection of degree $p^m$. The dual of that first isogeny is
\'etale, with unit differential multiplier; composition is $[p^m]$,
so the first multiplier has valuation $m$.
The remaining kernel maps isomorphically to its subgroup in the
\'etale quotient. To check this integrally, split that quotient after
a finite base extension. Each of its relevant components receives a
finite flat rank-one component of the remaining kernel, hence the
base itself. Thus the remaining isogeny is \`etale and its multiplier
is a unit. This proves $v_p(\lambda_C)=m$.
Equation \eqref{ell:ordinary-isogeny-difference} now gives
$v_p(h(E_C)-\upsilon_p(E_C))=b+n-2m$.
The integral second term proves the exact asserted formula for
$d_p(C)$: if this valuation is negative it cannot cancel with an
integral number, and otherwise both summands are integral.
Changing between good short N\'eron frames multiplies $h(E_C)$ by a
unit square, so the negative valuation is independent of that choice.

Choose a basis of $E[p^n]\simeq(\mathbf Z/p^n)^2$ with the
multiplicative subgroup as first axis. Primitive cyclic lines have
unique representatives $(a,1)$, with $a$ arbitrary, or $(1,z)$ with
$z\in p\mathbf Z/p^n$. The first set gives $m=0$ and has size $p^n$;
in the second, $m=v_p(z)$, with $z=0$ giving $m=n$.
This proves every stated branch count. As $b\ge0$, the sum of
$d_p(C)$ is at most
\[
 S_n=n+(p-1)\sum_{1\le k<n/2}p^{k-1}(n-2k).
\]
Subtracting the same expression for $n-2$ gives
$S_{2r}-S_{2r-2}=2p^{r-1}$ and
$S_{2r+1}-S_{2r-1}=(p+1)p^{r-1}$.
The initial values $S_1=1$, $S_2=2$ give the formulas above.

For non-CM $E$, $I_p$ is finite by the stronger adelic statement
\cite[Corollary 9.3, p.2393; restated p.2348]{EllLombardo2015OpenImage}.
Reduction modulo $p^n$ has index at most $I_p$ in a group acting
transitively on the $D_n$ cyclic lines. For a transitive finite
$G$-set and $H\subset G$ of index at most $I$, the orbit formula gives
\[
 |Hx|=\frac{|H|}{|H\cap\operatorname{Stab}_G(x)|}
 \ge\frac{|G|}{I|\operatorname{Stab}_G(x)|}=\frac{|Gx|}{I}.
\]
Thus $|\mathcal O|\ge D_n/I_p$. Nonnegativity bounds the sum over
this orbit by the full sum, proving the last assertion.
\end{proof}

The three exact programs
\texttt{check\_etale\_integral\_tangent.py},
\texttt{check\_isogeny\_hodge\_orbit.py}, and
\texttt{check\_isogeny\_branch\_denominators.py}
retain the local, degree-two, and cyclic-line certificates.
On the unique fully connected branch the negative valuation is
$\max(0,n-b)$; the proved orbit average tends to zero.
These are local N\'eron-frame norms. A global height statement would
also require a coherent metrized line and its other places; it is
not obtained by choosing unrelated global scalar frames.
The calculation supplies no contradiction to the nonzero
algebraic-$h$ case.
